% Generated draft: requires mathematical and source review.
\documentclass[11pt]{article}
\usepackage[margin=25mm]{geometry}
\usepackage{fontspec}
\setmainfont{Latin Modern Roman}
\usepackage{amsmath,amssymb,mathtools,unicode-math}
\setmathfont{Latin Modern Math}
\usepackage{enumitem,needspace,xurl,hyperref}
\setlength{\parindent}{0pt}
\setlength{\parskip}{5pt}
\setlength{\emergencystretch}{3em}
\newcommand{\N}{\mathbb{N}}
\newcommand{\Z}{\mathbb{Z}}
\newcommand{\Q}{\mathbb{Q}}
\newcommand{\R}{\mathbb{R}}
\newcommand{\C}{\mathbb{C}}
\newcommand{\F}{\mathbb{F}}
\newcommand{\A}{\mathbb{A}}
\newcommand{\PP}{\mathbb{P}}
\newcommand{\E}{\mathbb{E}}
\newcommand{\eps}{\varepsilon}
\newcommand{\dd}{\,\mathrm{d}}
\begin{document}
\section*{UnsolvedMath ID: 30001496}
\textbf{Source problem code:} catalogue entry
\subsection{Definitions and mathematical statement}
Let \(S\) denote the Enriques surface or its Fano model intended in the supplied record, and let \(D\) denote the divisor implicit in that model. The record does not specify the base field, the precise definition of the Fano model, or the conditions defining \(D\). Retain the supplied hypothesis that the complete linear system \(|2D|\), namely the system of effective divisors associated with \(\mathcal O_S(2D)\), gives an embedding
\[
\iota_{2D}\colon S\hookrightarrow\mathbb P^{19},
\qquad \deg\iota_{2D}(S)=40.
\]
The qualification that the Fano model is \emph{general} is part of the proposed assertion; the parameter space and exceptional locus are not supplied.

For clarity, the usual EPW construction uses a six-dimensional vector space \(V\), a volume form identifying \(\bigwedge^6V\) with the base field, and the resulting alternating pairing on \(\bigwedge^3V\) given by wedge product. For a ten-dimensional Lagrangian subspace \(A\subset\bigwedge^3V\), put
\[
F_v=v\wedge\bigwedge^2V,
\qquad
Y_A=\bigl\{[v]\in\mathbb P(V):A\cap F_v\ne\{0\}\bigr\}.
\]
Here the degeneracy locus has its determinantal scheme structure. When this construction gives a sextic hypersurface, \(Y_A\subset\mathbb P(V)\cong\mathbb P^5\) is an EPW sextic. Write \(\operatorname{Sing}(Y_A)\) for its singular locus. This terminology does not resolve the omitted base-field hypothesis or identify the particular EPW sextic intended by the source.

\textbf{Proposed statement, with unresolved qualifications.} Interpreting a suitable projection as a linear projection defined everywhere on \(\iota_{2D}(S)\), ask whether, for a general Fano model as above, there exist a projective linear subspace \(\Lambda\cong\mathbb P^{13}\subset\mathbb P^{19}\) and an EPW sextic \(Y_A\subset\mathbb P^5\) such that
\[
\Lambda\cap\iota_{2D}(S)=\varnothing,
\qquad
T=\pi_\Lambda\bigl(\iota_{2D}(S)\bigr),
\qquad
\deg T=40,
\qquad
T\subseteq\operatorname{Sing}(Y_A),
\]
where \(\pi_\Lambda\colon\mathbb P^{19}\dashrightarrow\mathbb P^5\) is projection from \(\Lambda\), and \(T\) is taken with its reduced image structure. The everywhere-defined projection and reduced-image conventions make this displayed version precise but are not explicitly specified in the record; the intended scheme-theoretic meaning of containment remains unresolved.

The supplied original wording additionally asserts that the projected surface is \emph{isomorphic to the corresponding degree-\(40\) surface} in the singular locus of \emph{the} EPW sextic. Thus, if that wording is retained, the intended conclusion also includes
\[
T\cong T_{\mathrm{corr}},
\qquad
T_{\mathrm{corr}}\subseteq\operatorname{Sing}(Y_A),
\qquad
\deg T_{\mathrm{corr}}=40.
\]
Neither \(T_{\mathrm{corr}}\) nor its correspondence with the Fano model is defined in the supplied record. Consequently, the existential containment question in the clean statement cannot be identified unconditionally with the original isomorphism assertion. The record does not assert that the projection restricts to an isomorphism from \(S\) onto \(T\).
\subsection{Short English statement}
The supplied conjecture asks whether a general Fano model of an Enriques surface, embedded by \(|2D|\) as a degree-\(40\) surface in \(\mathbb P^{19}\), admits a suitable projection to a degree-\(40\) surface in the singular locus of an EPW sextic. The original wording also requires an isomorphism with a particular corresponding surface whose construction is missing. The base field, polarization, meaning of generality, and projection conditions require clarification.
\subsection{Sources}
\begin{itemize}
\item Local UnsolvedMath registry: supplied complete source record for permanent ID \(30001496\), title \emph{EPW Sextics from Fano Models of Enriques Surfaces}, exact source problem code \texttt{catalogue entry}, in the Oberwolfach Reports Open Problems collection. No individual proposer is supplied; the proposed year is 2010.
\item Supplied primary reference: \emph{Classical Algebraic Geometry} (2010), Oberwolfach Reports, pp.\ 1573--1623, DOI \texttt{10.4171/owr/2010/27}. \url{https://doi.org/10.4171/owr/2010/27}. The literature-source list labels this the original OWR Report 27/2010.
\item Supplied related-literature reference, labeled ``Dolgachev--Markushevich related constructions'': \url{https://arxiv.org/abs/1906.01445}. No more detailed bibliographic information is supplied.
\item Literature assessment dated \textbf{2026-08-21}: the supplied curation labels the problem open, reports that the original report presents the general Fano-model assertion as a conjecture, and states that later related constructions do not establish this exact assertion. This is supplied dated curation, not an independently verified literature assessment.
\item The registry assigns the statement the label \texttt{corrected\_verified} and reports comparison with the official Oberwolfach report. That verification claim is supplied metadata; the sources have not been independently consulted for this draft.
\end{itemize}
\end{document}
